\section{Generación de comentarios}
\label{sec:generacion-comentarios}

Tras obtener las acciones del partido, el sistema ya dispone de la información necesaria para generar un comentario textual: la acción detectada, el instante en el que ocurre, el jugador asociado, su posición futbolística, su equipo y, cuando procede, información contextual adicional. El objetivo de esta fase no es construir una frase fija mediante reglas, sino utilizar un modelo generativo que permita obtener comentarios más naturales y con cierta variabilidad, manteniendo al mismo tiempo un control estricto sobre la información que puede aparecer en la respuesta.

\subsection{Variables de entrada del modelo}

La entrada al modelo se organiza como un evento estructurado en formato JSON. Esta decisión permite separar claramente dos partes del problema: por un lado, el pipeline anterior calcula la información objetiva del partido; por otro, el modelo de lenguaje transforma esa información en una frase natural de narrador.

Las variables principales utilizadas son las siguientes:

\begin{itemize}
    \item \texttt{action}: acción que se quiere comentar. Puede tomar valores como \texttt{control}, \texttt{pase}, \texttt{robo}, \texttt{tiro}, \texttt{corner}, \texttt{fuera de banda}, \texttt{saque de puerta}, \texttt{gol}, \texttt{intro} o \texttt{contexto}.

    \item \texttt{event\_time\_s}: segundo del vídeo en el que ocurre el evento.

    \item \texttt{player\_name}: nombre del jugador asociado a la acción.

    \item \texttt{player\_position}: posición futbolística del jugador, por ejemplo \texttt{MC}, \texttt{LI}, \texttt{ED} o \texttt{DC}.

    \item \texttt{team\_name}: equipo del jugador que realiza la acción.

    \item \texttt{opponent\_team\_name}: equipo rival. Se utiliza especialmente en goles e introducciones del partido.

    \item \texttt{team\_in\_favor}: equipo beneficiado en acciones como córner, saque de banda o saque de puerta.

    \item \texttt{action\_target}: destinatario u objetivo de la acción cuando existe esta información.

    \item \texttt{field\_zone}: zona del campo donde ocurre la acción, normalmente \texttt{zona de iniciación}, \texttt{zona de creación} o \texttt{zona de finalización}.

    \item \texttt{play\_context}: información contextual adicional. En comentarios de contexto puede incluir, por ejemplo, referencias tácticas a la defensa, al centro del campo o al ataque de un equipo.

    \item \texttt{match\_score}: contexto simulado de clasificación o situación del partido. Se utiliza para aportar información contextual, no para representar necesariamente el resultado real.

    \item \texttt{intensity}: marca ciertos eventos como interrupciones cuando aparece una acción peligrosa, como un tiro o un gol, durante un comentario de contexto.

    \item \texttt{action\_index}: índice secuencial de la acción dentro de la lista de eventos generada. Sirve principalmente para trazabilidad y depuración.
\end{itemize}

Estas variables permiten que el modelo disponga de información suficiente para generar comentarios coherentes sin inventar datos externos. Esta restricción es importante, ya que en una retransmisión automática resulta preferible generar un comentario sencillo pero correcto antes que una frase más vistosa con información no observada por el sistema.

\subsection{Diseño del prompt}

La generación del comentario se basa en dos prompts diferenciados: el \textit{system prompt} y el \textit{user prompt}. El primero define el rol general del modelo, mientras que el segundo contiene las instrucciones concretas para el evento actual.

El \textit{system prompt} establece el comportamiento global del modelo. Dependiendo del tipo de evento, el sistema cambia ligeramente el rol solicitado. Para una introducción, se le indica que actúe como un narrador de televisión abriendo la retransmisión antes de que comience el partido. Para un comentario de contexto, se le pide que actúe como comentarista de apoyo y que no narre una acción técnica concreta. En el caso de un gol, se permite un tono más expresivo, pero se imponen restricciones para evitar exageraciones o invenciones. Para el resto de acciones, se solicita una frase corta, natural y directa.

El \textit{user prompt}, en cambio, se construye de forma específica para cada evento. En él se incluyen reglas concretas como mencionar o no el minuto, utilizar literalmente la acción indicada, no inventar receptores en un pase, no añadir consecuencias posteriores a un gol o no repetir una frase reciente. Además, se adjunta el evento completo en formato JSON para que el modelo tenga acceso explícito a los datos disponibles.

Esta separación resulta útil porque permite mantener una personalidad general constante, pero ajustar las instrucciones según la naturaleza de cada evento. Por ejemplo, un comentario de introducción no debe mencionar acciones que todavía no han ocurrido, mientras que un gol debe mencionar claramente al autor, el equipo que marca y el equipo que encaja.

También se incorpora un mecanismo para reducir la repetición. Cuando el sistema dispone de un comentario reciente para una acción parecida, se añade al prompt una instrucción indicando que no debe repetirse la misma frase, ni el mismo verbo principal, ni la misma estructura. De esta forma se intenta mantener variabilidad sin perder fidelidad a los datos del evento.

\subsection{Selección del modelo de lenguaje}

Para generar los comentarios se probaron distintos modelos ligeros mediante Ollama. La prioridad era encontrar un equilibrio entre calidad lingüística, seguimiento de instrucciones y coste de inferencia. Dado que el comentario debe generarse a partir de eventos cortos y muy estructurados, no era necesario utilizar un modelo de gran tamaño, sino uno capaz de respetar bien el prompt y producir frases naturales en español.

\begin{table}[H]
    \centering
    \scriptsize
    \begin{tabular}{c|c}
        \textbf{Modelo} & \textbf{Tamaño} \\
        \hline\hline
        \texttt{gemma4:e2b} & 7.2 GB \\
        \texttt{qwen2:1.5b} & 934 MB \\
        \texttt{deepseek-r1:1.5b} & 1.1 GB \\
        \texttt{qwen3:1.7b} & 1.4 GB \\
        \texttt{lfm2.5-thinking:1.2b} & 731 MB \\
        \texttt{falcon3:1b} & 1.8 GB \\
        \texttt{granite4:1b} & 3.3 GB \\
        \texttt{gemma3:1b} & 815 MB \\
        \texttt{qwen3.5:2b} & 2.7 GB \\
        \texttt{qwen3:0.6b} & 522 MB \\
        \texttt{llama3.2:1b} & 1.3 GB \\
        \texttt{qwen3.5:0.8b} & 1.0 GB \\
        \texttt{tinyllama:1.1b} & 637 MB \\
        \texttt{qwen3.5:4b} & 3.4 GB \\
        \hline
    \end{tabular}
    \caption{Modelos probados para la generación de comentarios con Ollama.}
    \label{tab:modelos_llm_probados}
\end{table}


A partir de estas pruebas se observó que los modelos más pequeños de Qwen ofrecían una latencia favorable, pero la calidad de los comentarios no era suficiente para el tono buscado. En cambio, las variantes de mayor tamaño de Qwen mejoraban claramente la calidad lingüística, aunque su coste de inferencia aumentaba y dejaban de ser tan convenientes para una generación ágil dentro del sistema.

También se probó \texttt{lfm2.5-thinking:1.2b}, que generaba respuestas de buena calidad, pero presentaba una latencia demasiado alta para el objetivo del proyecto. Finalmente, el modelo que ofreció el mejor equilibrio observado entre naturalidad, seguimiento de instrucciones y tiempo de respuesta fue \texttt{gemma4:e2b}, por lo que se seleccionó como modelo principal para la generación de comentarios.

El principal inconveniente de utilizar directamente este modelo era su tamaño. Para reducir el consumo de recursos se optó por una versión cuantizada en formato \texttt{GGUF}, concretamente \texttt{gemma-4-E2B-it-Q4\_K\_S.gguf}. La cuantización consiste en representar los pesos del modelo con menor precisión numérica, reduciendo así el tamaño del modelo y el coste de inferencia. En este caso, \texttt{Q4} indica una cuantización principal a 4 bits, \texttt{K} hace referencia a una familia moderna de cuantizaciones por bloques, y \texttt{S} indica una variante más compacta.

Esta versión cuantizada ocupa aproximadamente 2.9 GB en disco, frente al tamaño mayor de la versión gestionada inicialmente mediante Ollama. Para ejecutarla se utilizó \texttt{llama.cpp}, una herramienta que permite desplegar modelos en formato \texttt{GGUF} de forma eficiente y exponerlos mediante una interfaz compatible con peticiones de chat. En la configuración del proyecto se utiliza un contexto de 512 tokens, suficiente para prompts cortos basados en eventos estructurados, y se desactiva el razonamiento explícito para evitar que el modelo genere explicaciones intermedias.

\subsection{Control de la salida generada}

Una vez recibido el comentario del modelo, el sistema aplica una limpieza básica de texto. Se eliminan espacios innecesarios, se normaliza la salida y se fuerza que el resultado final sea únicamente el comentario, sin explicaciones adicionales. Esto es importante porque algunos modelos tienden a responder con prefijos como ``Comentario:'' o a incluir aclaraciones que no forman parte de la frase final.

Además, el sistema evita comentar dos veces seguidas exactamente el mismo evento. Para ello se guarda una identidad simplificada del último evento generado, compuesta por la acción, el jugador y el equipo asociado. Si el siguiente evento coincide con esa misma identidad, se puede descartar para reducir duplicados.

En conjunto, esta fase transforma los eventos estructurados producidos por el pipeline en comentarios textuales naturales. La clave no está únicamente en elegir un modelo generativo, sino en controlar cuidadosamente qué información recibe, qué instrucciones debe seguir y qué restricciones impiden que invente contexto no observado por el sistema.

\subsection{Resultados de generación con eventos JSON}

Para comprobar el comportamiento del generador se levantó el modelo \texttt{gemma-4-E2B-it-Q4\_K\_S.gguf} mediante \texttt{llama.cpp}, utilizando el alias \texttt{gemma4-q4ks-text}. Las pruebas se realizaron con el servidor ya cargado en memoria.

Aunque la clase interna del evento define un conjunto amplio de campos posibles, el JSON que se introduce finalmente en el prompt no incluye siempre todas las claves: se eliminan los campos vacíos y se añade el campo \texttt{should\_mention\_minute}, que indica si el modelo debe mencionar explícitamente el minuto de partido.

\begin{table}[H]
    \centering
    \scriptsize
    \begin{tabular}{c|c|c|c}
        \textbf{Caso} & \textbf{Origen del evento} & \textbf{Tiempo total} & \textbf{Tiempo interno del modelo} \\
        \hline\hline
        Intro & Evento inicial con equipos & 0.311 s & 0.282 s \\
        Contexto & Evento construido desde alineación & 0.251 s & 0.238 s \\
        Pase & Evento real generado por PathCRF & 0.107 s & 0.094 s \\
        Tiro & Evento compatible con salida PathCRF & 0.101 s & 0.091 s \\
        Tiro con interrupción & Evento urgente con interrupción & 0.120 s & 0.113 s \\
        \hline
    \end{tabular}
    \caption{Tiempos de generación de comentarios con Gemma cuantizado ejecutado mediante \texttt{llama.cpp}.}
    \label{tab:tiempos_generacion_gemma}
\end{table}

\subsubsection*{Prueba 1: introducción}

\begin{verbatim}
{
  "action": "intro",
  "event_time_s": 0.0,
  "team_name": "Real Madrid",
  "opponent_team_name": "Wolfsburgo",
  "should_mention_minute": false
}
\end{verbatim}

Comentario generado:

\begin{quote}
¡Bienvenidos a todos! Ya está todo listo para este encuentro entre el Real Madrid y Wolfsburgo. ¡Que empiece la acción!
\end{quote}

\subsubsection*{Prueba 2: comentario de contexto}

\begin{verbatim}
{
  "action": "contexto",
  "event_time_s": 0.0,
  "team_name": "Real Madrid",
  "opponent_team_name": "Wolfsburgo",
  "field_zone": "zona de iniciacion",
  "play_context": "defensa de Real Madrid: centrales pepe y ramos",
  "match_score": "Real Madrid llega 3 en la tabla con 53 puntos; Wolfsburgo aparece 6 con 46 puntos",
  "action_index": 1,
  "should_mention_minute": false
}
\end{verbatim}

Comentario generado:

\begin{quote}
Con la defensa de Pepe y Ramos en el campo, el Real Madrid mantiene una ventaja en la clasificación, lo que les da tranquilidad en esta fase de salida de balón.
\end{quote}

\subsubsection*{Prueba 3: pase}

\begin{verbatim}
{
  "action": "pase",
  "event_time_s": 0.0,
  "player_name": "Carvajal",
  "player_position": "LD",
  "team_name": "Real Madrid",
  "action_target": "vini",
  "action_index": 1,
  "should_mention_minute": false
}
\end{verbatim}

Comentario generado:

\begin{quote}
Carvajal realiza un pase a Vini.
\end{quote}

\subsubsection*{Prueba 4: tiro}

\begin{verbatim}
{
  "action": "tiro",
  "event_time_s": 24.2,
  "player_name": "cristiano",
  "player_position": "DC",
  "team_name": "Real Madrid",
  "field_zone": "zona de finalizacion",
  "action_index": 28,
  "should_mention_minute": false
}
\end{verbatim}

Comentario generado:

\begin{quote}
Cristiano tira con potencia en la zona de finalización.
\end{quote}

\subsubsection*{Prueba 5: tiro con interrupción}

\begin{verbatim}
{
  "action": "tiro",
  "event_time_s": 24.6,
  "player_name": "cristiano",
  "player_position": "DC",
  "team_name": "Real Madrid",
  "field_zone": "zona de finalizacion",
  "play_context": "hay que cortar un comentario de contexto porque aparece una accion peligrosa",
  "intensity": "interrupcion",
  "action_index": 29,
  "should_mention_minute": false
}
\end{verbatim}

Comentario generado:

\begin{quote}
Perdon, te corto, Cristiano tira en la zona de finalización.
\end{quote}

Estos resultados muestran que el modelo responde de forma distinta según el tipo de evento. En acciones simples, como un pase o un tiro, genera frases breves y directas. En el evento de contexto utiliza la información táctica disponible y el contexto simulado de clasificación. Por último, cuando aparece \texttt{intensity = interrupcion}, el comentario adopta una forma de interrupción breve antes de narrar la acción peligrosa.

